% LSPV-5 REPORT::STRATFLAT1::2026-10-02T10:52:00-07:00
\documentclass[11pt]{article}
\usepackage[margin=.88in]{geometry}
\usepackage{amsmath,amssymb,booktabs}
\begin{document}
\section*{STRATFLAT1: polynomial-alternation replay collapses under a succinct source strategy}

\textbf{STATUS: PROVED/EXACT local flattening theorem + COMPUTATIONAL source control; independent red-team review pending. P versus NP remains OPEN.}

\section*{CORPUS REVIEWED}
The current GitHub tree was enumerated at 299 paths before the STRATFLAT1 writes. This run reviewed root \texttt{main.tex}; \texttt{AGENTS.md}; \texttt{CODEX.md}; the pnp-research skill and synthesis playbook; TASK-001 through TASK-007; verified MAIN snapshots through 1012; predecessor LSPV-4 TRACECLOSE1; LSPV-1 MATRIXCERT1 and current structural guards; LSPV-2 EXCESS1/EXCESSDUAL1/TURANEX1; LSPV-3 RAMCOMP1; ACR1, SBACR1, EVACR1; PRIME-LQP/PRIMECRT/PH-PROV-CLOSE; HEC/VSC-HEC/G-PPC-HEC/CRS-PPC/CCSPC/RCCN/Search--G; the current independent GROK review and hourly distill; governance; and cited ACR/product-span experiments and raw results.

The repository preflight script could not be executed locally because the active container does not contain a checkout of the repository. Its source/governance requirements were read through GitHub, and all mandatory GitHub-side checks are repeated after the writes.

Predecessor preflight is nonblocking: the newest LSPV-4 report is the substantive verified TRACECLOSE1 artifact with the expected sentinel.

\section*{DUPLICATION / CONTRADICTION AUDIT}
ACR1 already showed that polynomial authenticated replay does not generically escape TRACECLOSE1 or PH-PROV-CLOSE1. ACRTREE1 showed that an exponential explicit challenge tree can have a succinct strategy generator. SBACR1 made the game source-matched and found injective finite response signatures, but full reread/recomputation defeated reacquisition hardness. EVACR1 showed that compositional restrictions admit a succinct history quotient and therefore do not justify multiplicative cost.

STRATFLAT1 addresses the remaining syntactic hope directly: whether making the replay game have polynomially many alternating challenge/response rounds pushes the source-bound construction outside fixed PH. In the natural SU05/Hadamard replay model the answer is no, because the existential responses already have polynomial-time Skolem functions.

\section*{Primary-literature audit}
Chandra--Kozen--Stockmeyer establish \(AP=PSPACE\), so unrestricted polynomial alternation is a real complexity boundary. But QBF strategy literature emphasizes that winning moves are functions of prior universal choices and that these functions can have succinct circuit representations. Beyersdorff--Blinkhorn--Mahajan build partial QBF strategies succinctly into proofs. Heule--Seidl--Biere extract Skolem functions from QRAT proofs. Akshay--Beyersdorff--Chakraborty--Kasche--Mahajan--Spachmann (SAT 2026) study proof systems from which Skolem and Herbrand functions and intermediate game values can be extracted efficiently. These are direct precedents for the strategy-representation issue, but they do not transfer QBF proof lower bounds to unrestricted SAT.

Santhanam--Williams connect faster algorithms for quantified Boolean formulas to circuit lower bounds. This reinforces that QBF/strategy complexity is its own load-bearing problem rather than a free consequence of SAT provenance.

\section*{Exact game}
Let \(T(n)\) be polynomially bounded. A source-bound alternating replay game has a polynomial-length committed source \(R\), universal challenges \(c_i\), existential openings \(o_i\), and a polynomial-time verifier \(V\):
\[
\exists R\;\forall c_1\exists o_1\;\forall c_2\exists o_2\;\cdots\;\forall c_T\exists o_T\;V(x,R,c_{\le T},o_{\le T}).
\]
The intended hope was that unbounded \(T\) makes the SAME\_REALIZATION requirement genuinely alternating and therefore escapes PH-PROV-CLOSE1.

\section*{PROVED/EXACT: deterministic-Skolem flattening}
Suppose there are polynomial-time functions \(F_i(x,R,c_{\le i})\) such that for every legal source and challenge prefix the required opening is \(o_i=F_i(x,R,c_{\le i})\), and substitution of these openings preserves the declared legality condition. Then the alternating condition is equivalent to
\[
\boxed{\exists R\;\forall C\;W(x,R,C)}
\]
where \(C=(c_1,\ldots,c_T)\) is one polynomial-length challenge string and
\[
W(x,R,C)=V\bigl(x,R,C,F_1(x,R,c_1),\ldots,F_T(x,R,c_{\le T})\bigr).
\]
Because \(T\) is polynomial and every \(F_i,V\) is polynomial-time, \(W\in P\). Therefore the source-existence language is in \(\Sigma_2^P\), despite the displayed polynomial number of alternations.

\paragraph{Proof.}
For fixed \(x,R,C\), each existential opening is uniquely supplied by its stated polynomial-time function. Eliminating the existentially quantified openings by substitution leaves only the one universal quantifier over the concatenated challenge sequence. Conversely, the substituted openings witness every existential move in the original game. Polynomially many polynomial-time evaluations remain polynomial-time. \(\square\)

Under the counterfactual \(P=NP\), \(\Sigma_2^P=P\), so a polynomially balanced valid source is again selectable in FP by PH-PROV-CLOSE1.

\section*{PROVED/EXACT: succinct-strategy flattening}
More generally, suppose a winning strategy has a polynomial-size description \(S\) and a polynomial-time evaluator \(\operatorname{Eval}(S,i,c_{\le i})=o_i\). Then existence of such a succinct winning strategy is expressible as
\[
\boxed{\exists R,S\;\forall C\;V\bigl(x,R,C,\operatorname{Eval}(S,C)\bigr),}
\]
again a \(\Sigma_2^P\) predicate when \(|R|,|S|\) are polynomially bounded.

Thus polynomial alternation escapes the fixed-PH closure only when the proof needs a strategy for which no polynomial-size/evaluable representation is available. Proving that all winning strategies require superpolynomial general circuits would itself be a substantive strategy/circuit lower bound. No such source-matched theorem exists in this project.

\section*{Exact SU05/Hadamard consequence}
In SBACR1, the committed source is the Boolean-grid word \(w\), a challenge is an orbit shift \(s\), and the response is a basis for the legal slice \(S_s(w)\). The canonical response
\[
F(w,s)=\operatorname{RREFBasis}(S_s(w))
\]
is polynomial-time computable from the static source by interpolation, orbit evaluation, binary labeling and Gaussian elimination. It is independent of prior response choices.

Therefore even the apparently polynomial-alternating game
\[
\exists w\;\forall s_1\exists B_1\;\cdots\;\forall s_T\exists B_T:\bigwedge_i [B_i\text{ spans }S_{s_i}(w)]
\]
has a fixed polynomial-time Skolem strategy \(B_i=F(w,s_i)\). The alternation is syntactic rather than resource-forcing.

\section*{COMPUTATIONAL control STRATFLAT1}
Artifacts:
\[
\texttt{experiments/acr/strategy\_flattening.py},\qquad\texttt{results/raw/acr/2026-10-02\_LSPV5\_STRATFLAT1.json}.
\]
The run is exhaustive and deterministic.

For \((p,h,m,T)=(17,4,4,4)\), all 14 nonconstant sources and all \(16^4=65{,}536\) challenge sequences per source were checked, for 917,504 source-paths total. Interactive canonical responses exactly matched flattened direct evaluation in every case. Of 224 source-slices, 220 (98.2\%) had rank at least two and therefore at least two distinct legal ordered-basis encodings by basis-row permutation.

For \((29,5,5,3)\), all 30 nonconstant sources and all \(28^3=21{,}952\) challenge sequences per source were checked, for 658,560 source-paths. Again every interactive response matched the flattened strategy. Of 840 source-slices, 839 (99.88\%) admitted multiple ordered-basis encodings.

Thus local response nonuniqueness is not enough. Even when most challenges have several legal existential answers, canonical RREF gives one polynomial-time winning selector.

\section*{Static-SAT compiler attack}
The hoped-for compiler would have to take every hypothetical polynomial-time static SAT decider \(A\) and force it into a source-bound game for which polynomial alternation creates an unavailable resource. STRATFLAT1 splits this into two cases.

\paragraph{Case 1: polynomial source/strategy representation.}
If the compiled SAME\_REALIZATION has a polynomial-size source or strategy generator and legal openings are polynomially evaluable, the game flattens to fixed PH. Under \(P=NP\) it is inside the existing closure ladder. This includes the current SBACR1 source and EVACR1 compositional restriction quotient.

\paragraph{Case 2: no polynomial strategy representation.}
Then the compiler must prove that the induced source-legal game has no polynomial-size/evaluable winning strategy. That is a new representation/circuit lower bound on the strategy functions, and it must hold for the object forced by every arbitrary static SAT decider. Merely defining a QBF, DQBF, PSPACE game, or exponentially large strategy tree does not prove this transfer.

Hence the proposed static-SAT-to-polynomial-alternation compiler is
\[
\boxed{\textbf{FALSIFIED/RETIRED in its natural source-computable form}.}
\]
The only repair is to supply a source-matched superpolynomial strategy-representation lower bound plus an arbitrary-algorithm necessity theorem.

\section*{Nonconfluence and EVACR1}
EVACR1 already shows that ordinary restrictions compose to a compact history quotient. Nonconfluent evolution can make existential choices strategically relevant, but STRATFLAT1 adds a sharper guard: nonconfluence matters only if it destroys every polynomial-size winning strategy representation. If a polynomial circuit/Skolem map still evaluates the strategy, the polynomial alternation again flattens to \(\Sigma_2^P\).

This identifies the true missing theorem as strategy-representation hardness, not alternation count.

\section*{Barrier and hostile-control audit}
\begin{itemize}
\item \textbf{Relativization:} the flattening theorem relativizes. It is a closure theorem, not a nonrelativizing hinge.
\item \textbf{Algebrization:} the current SU05 source evaluator is algebraic/interpolative, so algebraic oracle access only strengthens the recomputation attack.
\item \textbf{DEAD/COPY/DELAY:} syntactic strategy histories can be inflated without changing the canonical source evaluator.
\item \textbf{INTERP/RECOMP:} polynomial universal simulation or recomputation preserves succinct strategy access.
\item \textbf{Easy controls:} Horn, 2-SAT, affine/Tseitin, bounded-treewidth, padding, renaming and independent products also admit succinct strategies when encoded in this replay format; alternation count does not distinguish them.
\item \textbf{QBF proof complexity:} lower bounds for restricted QBF proof systems or restricted strategy representations are FAILED TRANSFERS to unrestricted SAT absent a compiler.
\item \textbf{DQBF:} allowing arbitrary dependency Skolem functions can raise the ambient problem well beyond QBF, but importing DQBF/NEXP-level synthesis into the definition is not a SAT lower bound.
\end{itemize}

\section*{CORPUS RECONCILIATION}
\textbf{New verified local theorem:} deterministic-Skolem and polynomial-succinct-strategy flattening.

\textbf{Conflict resolved:} polynomially many displayed challenge/response quantifiers do not by themselves leave fixed PH when existential responses have a succinct strategy representation.

\textbf{Stale claim removed/downgraded:} the suggestion that polynomial alternation alone is the next provenance escape is too weak. The required escape is polynomial alternation plus a source-matched lower bound excluding all polynomial-size/evaluable winning strategies.

\textbf{Falsifications retained:} explicit strategy-tree size, response-signature diversity alone, prime provenance, deterministic trace provenance, canonical binding, ordinary evolving restrictions, and multiplicative cost from branch count remain retired as generic hinges.

\textbf{Adopted handoffs:} ACR1, SBACR1, EVACR1 as hostile controls; PH-PROV-CLOSE1; TRACECLOSE1; current structural BITRANK/MATRIXCERT/EXCESS/TURANEX branch.

\section*{WHAT CHANGED RELATIVE TO ROOT MAIN}
The provenance branch is narrowed again. The 1012 target's source-bound adaptive SAME\_REALIZATION phrase remains useful only as an instrumentation problem. Polynomial alternation is not itself a surviving acquisition mechanism. A provenance proof-spine candidate now requires an explicit lower bound on the representation/evaluation complexity of every legal winning strategy, plus a static-SAT compiler forcing that game.

The structural characteristic-two target remains the highest-value theorem target.

\section*{Single next target}
Do not add more challenge rounds. On the exact source-typed SU/vMS family, ask whether the response map required by the structural reconstruction can be proved to have superpolynomial circuit/Skolem complexity for the same arbitrary solver realization. First run the hostile opposite test: search for compact circuits, algebraic evaluators, decision diagrams, or source quotients computing all legal responses. If any polynomial representation survives, provenance remains instrumentation only.

\begin{thebibliography}{9}
\bibitem{CKS81} A. K. Chandra, D. C. Kozen, L. J. Stockmeyer, Alternation, J. ACM 28(1), 1981, 114--133.
\bibitem{BBM21} O. Beyersdorff, J. Blinkhorn, M. Mahajan, Building Strategies into QBF Proofs, J. Automated Reasoning 65, 2021.
\bibitem{HSB14} M. Heule, M. Seidl, A. Biere, Efficient Extraction of Skolem Functions from QRAT Proofs, FMCAD 2014.
\bibitem{Akshay26} S. Akshay, O. Beyersdorff, S. Chakraborty, L. Kasche, M. Mahajan, L. N. Spachmann, Proof Systems for QBF Synthesis: Extracting Skolem and Herbrand Functions, SAT 2026.
\bibitem{SW15} R. Santhanam, R. Williams, Beating Exhaustive Search for Quantified Boolean Formulas and Connections to Circuit Complexity, SODA 2015.
\end{thebibliography}

\end{document}
% END_OF_VERIFIED_REPORT::LSPV-5::STRATFLAT1
